% Round-loop per-process phase machine (PLTS.ABA.Composition.CoreProcStepN). Automata style.
\begin{tikzpicture}[
    ->, >={Stealth[round]}, auto, semithick, node distance=2.9cm, on grid,
    every state/.style={draw, ellipse, minimum width=1.8cm, minimum height=0.9cm,
      inner sep=1pt, align=center, font=\small},
    lbl/.style={font=\footnotesize, fill=white, inner sep=1.5pt},
    elbl/.style={font=\scriptsize, fill=white, inner sep=1.2pt}]
  \node[state, initial, initial text={}] (idle) {idle};
  \node[state] (tcg) [right=of idle] {\texttt{toCallG}};
  \node[state] (awg) [right=of tcg]  {\texttt{awaitG}};
  \node[state] (tcw) [right=of awg]  {\texttt{toCallW}};
  \node[state] (aww) [right=of tcw]  {\texttt{awaitW}};
  \path
    (idle) edge node[elbl, above]{\texttt{input}} (tcg)
    (tcg)  edge node[elbl, above]{\texttt{callG}} (awg)
    (awg)  edge node[elbl, above]{\texttt{retG}} (tcw)
    (tcw)  edge node[elbl, above]{\texttt{callW}} (aww)
    % inputLoop sits at every phase but idle; it is drawn at toCallG, on the
    % upper left so that the round back-edge can come down on the upper right
    (tcg)  edge[loop, out=110, in=150, looseness=7]
      node[lbl, above left, text width=1.6cm, align=center]{\texttt{inputLoop}} (tcg);
  % round back-edge, routed over the top so it clears the DECIDED node below
  \draw[->] (aww) to[out=115, in=50, looseness=0.55]
    node[lbl, above]{\texttt{retW}: round${+}1$ (D10)} (tcg);
  % the DECIDED sets, below the row; detail in caption
  \node[draw, dashed, rounded corners, align=center, font=\footnotesize,
        below=1.95cm of awg, minimum width=5.4cm, inner sep=4pt] (dec)
    {\textbf{the DECIDED sets}\\ ($\tau$-loops, D12$'$)};
  \draw[dotted, ->] (dec.north) -- (aww.south);
\end{tikzpicture}
